% ============================================================
% Einleitung -- trägt die Kapitelnummer 1 (Handbuch S. 83)
%
% Inhalt gemäss Handbuch:
% - Untersuchungsgegenstand, Problemstellung, Theorie und Ziel der
%   Arbeit definieren
% - Untersuchungsfeld klar eingrenzen, Ausgangslage darstellen
%   (grösserer Zusammenhang, vorhandene Literatur, bestehende Versuche)
% - Am Schluss der Einleitung: Hypothesen, die sich aus dem Theorieteil
%   ableiten lassen
%
% Für dieses Thema z.B.: Warum Bier selber brauen? Welche Fragestellung
% wird untersucht (z.B. Einfluss von Gärtemperatur/Hefesorte auf
% Geschmack/Alkoholgehalt)? Welche Hypothese wird geprüft?
% ============================================================

\chapter{Einleitung}
\label{kap:einleitung}

% Beispieltext -- als Ausgangspunkt gedacht, noch zu ergänzen und mit
% eigenen Formulierungen zu ersetzen.

Bier zählt zu den ältesten industriell hergestellten Lebensmitteln der
Menschheit und wird bis heute nach einem im Kern unveränderten Prinzip
gebraut: Aus Wasser, Malz, Hopfen und Hefe entsteht durch einen
kontrollierten biochemischen Prozess ein alkoholhaltiges Getränk
\parencite{bruecklmeier2022}. Obwohl der Brauprozess auf den ersten
Blick einfach erscheint, hängt das Ergebnis von einer Vielzahl
kontrollierbarer Grössen ab, etwa der Maischetemperatur, der
Hopfenmenge oder der verwendeten Hefesorte.

In der vorliegenden Arbeit brauen wir in einer Gruppe von zwei Personen
ein eigenes Bier und dokumentieren den gesamten Prozess von der
Zutatenbestellung über den Brauvorgang bis zur Abfüllung in Flaschen.
Im Zentrum steht die Frage, wie sich die Gärtemperatur auf den
Vergärungsgrad und den Geschmack des fertigen Biers auswirkt. Wir
vermuten, dass eine höhere Gärtemperatur zu einem höheren
Vergärungsgrad, aber auch zu stärker ausgeprägten Fruchtnoten führt
\vgl[S.~45]{bruecklmeier2022}.

